\documentclass[11pt]{article}
\usepackage[margin=1in]{geometry}
\usepackage{fontspec}
\setmainfont{DejaVu Serif}
\setmonofont{DejaVu Sans Mono}
\usepackage{amsmath,amssymb}
\usepackage{fvextra}
\usepackage[unicode,colorlinks=true,linkcolor=blue,urlcolor=blue]{hyperref}
\tolerance=5000
\setlength{\emergencystretch}{4em}
\setlength{\parskip}{0.4em}
\title{P15.3: dense null leaves, valuation collapse, and the unresolved general construction}
\author{Unverified research draft}
\date{}
\begin{document}
\maketitle


This is an unverified mathematical draft. The new result is a counterexample to a precisely specified unrescaled current-\allowbreak{}valuation recipe, even for the ordinary smooth P-\allowbreak{}type. It is not an impossibility theorem for original P15.3. The supplied divisorial construction, transcendental product example, and dense-\allowbreak{}pole obstruction are retained in full in the appendix, with the two stipulated wording corrections. No independent correctness review of the new argument has been performed.

\subsection{1. Exact result and conventions}

Use $dd^c=i\partial\bar\partial/(2\pi)$, so that $dd^c\log|z|^2$ is the unit divisor current. A potential is the canonical upper semicontinuous quasi-\allowbreak{}psh representative and is not identically $-\infty$. For a smooth real closed $(1,1)$-\allowbreak{}form $\theta$, write

\[P_\theta[u]=\left(\sup\{v\in\operatorname{PSH}(X,\theta):v\leq0,\ v\leq u+C_v\}\right)^*.\]

This ordinary envelope is not an I-\allowbreak{}envelope. The prescribed datum in the new example is $u=0$, with $P_\beta[0]=0$. Its smooth bounded singularity type is retained without substituting multiplier ideals.

\textbf{Theorem.} There exist a compact Kähler complex surface $M$, with no curves, and a nonzero constant semipositive rank-\allowbreak{}one form $\beta$ on $M$ such that, for $\alpha=[\beta]$:

\begin{enumerate}
\item[1.] $\alpha$ is genuinely transcendental, nef, nonbig, and has zero class volume. Its nef intersection numerical dimension is one. The largest nonzero degree of the actual nonpluripolar products of $\beta$ and of Cao's current products is also one; those assertions are computed separately.
\item[2.] $\operatorname{PSH}(M,\beta)=\mathbb R$. For every proper bimeromorphic morphism $\pi:X\to M$ with $X$ a smooth compact Kähler surface, $\operatorname{PSH}(X,\pi^*\beta)=\mathbb R$. Thus $\beta$ and $\pi^*\beta$ are the unique positive closed currents in their respective classes.
\item[3.] Fix any such $\pi$, any globally smooth connected complete flag $X\supset Y\supset\{q\}$, and any Kähler form $\Omega$ on $X$. For the usual DRWNXZ current-\allowbreak{}valuation bodies,

\end{enumerate}
\[\lim_{\varepsilon\downarrow0}\Delta_{Y,q}(\pi^*\alpha+\varepsilon[\Omega])=\{(0,0)\}\]

in Hausdorff distance, and their intersection over all positive $\varepsilon$ is that point. The partial bodies with reference $\pi^*\beta+\varepsilon\Omega$ are the same bodies. The limit is also the same along any sequence of smooth closed Kähler forms $\rho_j\to0$ in uniform norm, in place of $\varepsilon\Omega$.

The quantification in (3) is over every fixed modification and every fixed smooth flag on it. It asserts no estimate uniform over a varying family of modifications or flags. The standard coordinate carrier is $\mathbb R^2$; the limiting affine carrier is the single point. There is no one-\allowbreak{}dimensional positive-\allowbreak{}volume complex subvariety on $M$, or for the pullback class on any of these models. This is exactly the obstruction to inferring the desired one-\allowbreak{}dimensional numerical carrier from the usual limiting body.

\subsection{2. Published inputs and classical background, with their actual domains}

Only the four permitted primary caches have been read. Their local byte-\allowbreak{}identical copies, headers, literal statements and proof bodies, and inclusive LF coordinates are in \texttt{source-\allowbreak{}extracts/\allowbreak{}source-\allowbreak{}provenance.json}. Coordinates always refer to one-\allowbreak{}based entries of \texttt{text.split('\textbackslash{}n')}; form feeds and the last empty entry are retained.

DRWNXZ is arXiv:2309.07584v3, 5 January 2026, 2034 LF entries. Its initial standing assumption is a big class on a connected compact Kähler manifold (LF 226–263). Its globally smooth connected complete flag and recursive generic-\allowbreak{}Lelong valuation are LF 561–595. Lemma 3.1, including its regularization/\allowbreak{}Briançon–Skoda proof, and Definition 3.2 are LF 597–620: the big-\allowbreak{}class body is the Euclidean closure of gentle analytic Kähler-\allowbreak{}current valuations, equivalently of positive-\allowbreak{}current valuations analytic near the flag point. LF 624–691 gives scaling, addition, continuity from above, and the pseudoeffective intersection definition. This proof uses those results at positive levels, which are Kähler here, and uses the intersection definition at zero.

Divisorial subtraction is Lemma 2.5, LF 314–333: for rational $a\geq0$, generic Lelong number at least $a$ implies positivity after subtracting $a[D]$, and local analytic singularities persist. The proof cites Siu decomposition for positivity and divides local ideal generators by the divisor equation for the analytic assertion. We retain that published positivity theorem. All coefficients subtracted from generating gentle currents below are rational, as in this source. Gentleness, rational analytic coefficients, and the source's restricted use of the word modification are LF 338–385. Our theorem quantifies over proper bimeromorphic morphisms with smooth compact Kähler source; its descent proof does not need their factorization into smooth blowups.

Definition 5.1/\allowbreak{}Example 5.3 and their proofs are LF 852–894. The partial body has a \textbf{gentle analytic Kähler} reference and imposes more singularity up to an input-\allowbreak{}dependent bounded constant. Smooth references at positive levels satisfy those hypotheses. The source's Theorem 7.6 and Corollary 7.7, LF 1389–1415, are class-\allowbreak{}level ambient volume results. Their final argument bounds ambient volume by perturbed big-\allowbreak{}class volumes. Their induction proof is retained in the extracts, including LF 939–963 and 1139–1388. No lower-\allowbreak{}dimensional numerical identity or prescribed arbitrary P-\allowbreak{}potential theorem is asserted there.

Xia is arXiv:2112.04290v2, 13 April 2024, 2700 LF entries. Equation (1.2), LF 67–105, uses sections of powers of a line bundle and multiplier ideals on a projective variety. LF 371–429 distinguishes bounded singularity comparison, the ordinary P-\allowbreak{}envelope, the mass-\allowbreak{}sensitive C-\allowbreak{}envelope, and the I-\allowbreak{}envelope. Proposition 2.10 identifies C with ordinary P only under positive \textbf{actual input top mass}; at zero mass its displayed limiting identity is different. Theorem 2.19, LF 476–494, is stated as a cited theorem and computes section volume from the actual top mass of the \textbf{I-\allowbreak{}envelope}, not unconditionally from the input potential. LF 1158–1185 records I-\allowbreak{}invariance and its proof reference. Section 5.6, LF 1623–1637, defines the rational ample-\allowbreak{}perturbation intersection and explicitly labels the dimension equality as Conjecture 5.29. This nonprojective counterexample is not a counterexample to that projective conjecture. The Chebyshev question, LF 1824–1840, concerns the auxiliary \textbf{more-\allowbreak{}singular Kähler-\allowbreak{}current potential $\eta$}; it is not a question about an arbitrary positive form.

Cao is arXiv:1210.5692v1, 21 October 2012, 1931 LF entries. LF 162–166 explicitly states the global potential representation of a closed positive $(1,1)$-\allowbreak{}current on a compact Kähler manifold. Definitions 2–3, Proposition 2.2 and its comparison proof, LF 172–279, concern quasi-\allowbreak{}equisingular analytic approximations and their absolutely continuous wedge \textbf{cohomology} products. Definition 4, LF 281–289, defines numerical dimension for a pseudoeffective \textbf{line-\allowbreak{}bundle metric}. We use Definition 3 for smooth transcendental currents, not Definition 4 as a line-\allowbreak{}bundle theorem about this torus. Its projective general-\allowbreak{}hypersurface restriction argument, LF 660–721, cannot supply a curve on this curve-\allowbreak{}free manifold.

CHPW is arXiv:1508.03922v5, header 5 February 2018 and internal date 24 September 2018, 1425 LF entries. Theorem B is LF 166–185. Definition 2.19 and the ample-\allowbreak{}divisor construction are LF 530–577. Theorem 3.20, Proposition 3.21, Corollary 3.22 and their proofs are LF 915–1085; Remark 3.24 is LF 1090–1096. The numerical equality requires a smooth projective variety, a pseudoeffective real \textbf{divisor}, and a flag containing a positive volume subvariety of the requisite numerical dimension. Its proof uses restricted graded series and then the flag slice/\allowbreak{}volume growth argument. Those hypotheses fail here. We do not apply that theorem to $\alpha$.

For completeness, the precise classical analytic background retained below is as follows. A compact Kähler manifold admits the $dd^c$ potential representation of positive currents just stated by Cao; that representation includes a canonical quasi-\allowbreak{}psh potential. A locally bounded-\allowbreak{}above psh function extends across an analytic subset by its upper semicontinuous extension, and two canonical psh representatives equal almost everywhere are equal everywhere; these are exactly the inputs used in DRWNXZ Lemma 2.1's proof, LF 255–263. We reuse these classical statements, while proving the extension of that argument to a nonbig class. Contraction of model curves is proved directly by pushing forward their integral integration currents; it requires no separate proper-\allowbreak{}image theorem. Weak compactness of bounded positive currents, the needed Liouville assertion, the torus cohomology calculation, and dense-\allowbreak{}leaf existence are reconstructed below. Ordinary smooth products need none of the general nonpluripolar continuity or mass-\allowbreak{}monotonicity results. The appendix's bounded-\allowbreak{}potential products use the classical Bedford–Taylor/\allowbreak{}nonpluripolar rules explicitly stated there; no new theorem about arbitrary unbounded potentials is imported from them.

\subsection{3. Constructing the curve-\allowbreak{}free torus and a dense null line}

Start with $V_{\mathbb R}=\mathbb R^4$ and $\Lambda=\mathbb Z^4$. Vary the linear complex structure $J$ in a fixed component of the space of oriented complex structures. This is a connected real-\allowbreak{}analytic manifold and a Baire space. Its constant forms compute the cohomology of the real torus: integration over coordinate two-\allowbreak{}tori identifies $H^2(V_{\mathbb R}/\Lambda,\mathbb Z)$ with integral alternating bilinear forms on $\Lambda$. One way to see the constant representative assertion is to average a closed form over all translations. Each translation is isotopic to the identity, so averaging preserves the cohomology class; a translation-\allowbreak{}invariant form is constant, and its coefficients are fixed by these periods. Averaging preserves type $(1,1)$ for a fixed $J$. Thus an integral class has type $(1,1)$ precisely when its constant alternating representative $e$ satisfies

\[e(Jv,Jw)=e(v,w)\quad\text{for all }v,w.\]

For each nonzero integral $e$, this is a closed real-\allowbreak{}analytic condition with empty interior. Here is the necessary properness check. Choose $v_1,v_3$ with $e(v_1,v_3)\ne0$, and complete to an oriented real basis $v_1,v_2,v_3,v_4$ for the chosen complex-\allowbreak{}orientation component. Rescale $v_2$ by a positive factor, if necessary, so that $e(v_2,v_4)\ne e(v_1,v_3)$. Define $Jv_1=v_2$, $Jv_2=-v_1$, $Jv_3=v_4$, $Jv_4=-v_3$ using this rescaled basis. The displayed invariance fails. This construction works in either orientation component. The identity principle for real-\allowbreak{}analytic functions on the connected component therefore makes the locus nowhere dense. There are only countably many nonzero integral $e$. Baire's theorem gives a $J$ for which no such class is of type $(1,1)$.

Put $V=(V_{\mathbb R},J)$ and $M=V/\Lambda$. Any positive constant $J$-\allowbreak{}Hermitian metric defines a constant closed Kähler form $\omega_M$ on $M$. A curve on this complex surface would be a nonzero effective divisor. Its integral $(1,1)$-\allowbreak{}class has pairing $\int_C\omega_M>0$, by integration on its regular locus (or its normalization), and is consequently nonzero. That contradicts the choice of $J$. Thus $M$ has no curves and the real span of divisor classes is zero.

Choose a complex line $W\subset V$ with dense image in the real torus. Existence can be checked without assuming a holomorphic quotient. For each nonzero character $k\in\Lambda^*=\mathbb Z^4$, the condition $k|_W=0$ removes exactly one line from $\mathbb P(V)\simeq\mathbb P^1$: that line is the complex kernel of $v\mapsto k(v)-i k(Jv)$. Choose $W$ outside those countably many lines.

Here is a direct density proof. Let $w,Jw$ span its real plane. Average a torus character $x\mapsto\exp(2\pi i k(x))$ over $x+sw+tJw$, with $s,t\in[-R,R]$, using normalized area. The average is a product of two sinc factors, times the initial character value. For nonzero $k$, at least one of $k(w),k(Jw)$ is nonzero, so that average tends to zero. The zero character averages to one. Trigonometric polynomial approximation then gives convergence of these averages for every continuous function to its torus Haar integral. A nonnegative continuous bump supported in an arbitrary nonempty open set has positive integral; its average is eventually positive. Thus every affine leaf $x+W$ enters every nonempty open set, as required. This is a mathematical Fourier argument, not a finite numerical test.

Let $\ell:V\to\mathbb C$ be a nonzero complex-\allowbreak{}linear functional with kernel $W$, and take a positive constant multiple of $i\,d\ell\wedge d\overline\ell$ as $\beta$. It descends to $M$, is closed and semipositive, has rank one and kernel $W$, and satisfies

\[\beta^2=0,\qquad A_M:=\int_M\beta\wedge\omega_M>0.\]

The second statement holds pointwise from its one positive Hermitian eigenvalue. It proves $[\beta]\ne0$, and hence genuine transcendence since the divisor span is zero.

\subsection{4. Rigidity of potentials and descent on every smooth Kähler model}

We first verify the Liouville assertion used for the null leaves. If $h$ is a nontrivial subharmonic function on $\mathbb C$, bounded above, choose an origin where it is finite. Jensen's formula expresses its circular mean as its value at that origin plus the integral of the positive Laplacian measure in discs, divided by their radii. If any disc has positive Laplacian mass, that mean grows at least as a positive constant times $\log R$ for large $R$, contradicting the upper bound. Hence the Laplacian is zero and $h$ is harmonic. For a positive harmonic function $B-h$ on a disc of radius $R$, the differentiated Poisson formula bounds its derivative at the center by $2(B-h(0))/R$. Letting $R\to\infty$ at any center makes all derivatives zero. Thus $h$ is constant. This also covers subharmonic functions with occasional $-\infty$ values, since zero Laplacian gives the canonical harmonic representative.

Take $\varphi\in\operatorname{PSH}(M,\beta)$. It is bounded above by compactness and has a finite value at some $x$. Along any affine holomorphic null leaf, its restriction is subharmonic or identically $-\infty$: locally add a smooth potential for $\beta$, whose restriction to that leaf is harmonic, since the restricted form is zero. The leaf through $x$ has a nontrivial restriction, so the Liouville argument makes its value the constant $c=\varphi(x)$. It is dense. For any $y$, take points on that leaf converging to $y$. Upper semicontinuity gives $\varphi(y)\geq c$, so $\varphi$ is finite everywhere. Now each leaf has a constant restriction. If two such constants are $c_1,c_2$, density of the first leaf and upper semicontinuity at a point of the second give $c_2\geq c_1$. The reverse argument gives $c_1\geq c_2$. Therefore $\varphi$ is globally constant. Constants are conversely $\beta$-\allowbreak{}psh. The global potential representation of positive currents then proves uniqueness of the positive current $\beta$ in its class.

Next let $\pi:X\to M$ be any proper bimeromorphic morphism with smooth compact Kähler source, and put $\beta_X=\pi^*\beta$. The nonbig descent proof is included rather than invoking a standing big-\allowbreak{}class result outside its domain. If $v\in\operatorname{PSH}(X,\beta_X)$, on the open set $U$ where $\pi$ is an isomorphism put $w=v\circ\pi^{-1}$. It is $\beta$-\allowbreak{}psh there and bounded above, because $v$ is bounded above on compact $X$. In a local chart add a smooth potential for $\beta$. The removable-\allowbreak{}singularity statement extends the resulting locally bounded-\allowbreak{}above psh function across the analytic complement of $U$. These extensions glue to a $\beta$-\allowbreak{}psh function $w$ on $M$. Its pullback and $v$ agree off an analytic subset, hence almost everywhere, and then everywhere as canonical quasi-\allowbreak{}psh representatives. Pullback is injective and reflects pointwise order since $\pi$ is surjective. This proves the order isomorphism of potential spaces without any bigness assumption. The torus rigidity makes $w$, hence $v$, constant. The same global potential representation proves uniqueness of $\beta_X$ in $\pi^*\alpha$.

Consequently the ordinary envelopes of the specified zero potential are exactly zero both downstairs and upstairs. There is no identification with an I-\allowbreak{}envelope or a claim that a cohomology product determines arbitrary P-\allowbreak{}types.

\subsection{5. Numerical data and the missing complex carrier}

The smooth semipositive form $\beta$ proves nefness: $\beta+\varepsilon\omega_M$ is Kähler for every positive $\varepsilon$. A hypothetical Kähler current $T$ in $\alpha$ with $T\geq\delta\omega_M$ would satisfy

\[0=\int_M T\wedge\beta\geq\delta\int_M\omega_M\wedge\beta>0.\]

The equality is cohomological and uses closedness; positivity is valid because the test form is smooth semipositive. Thus the class is nonbig, independently of the mass of any chosen potential. Its class volume is zero, as in the nonbig pseudoeffective convention of DRWNXZ Corollary 7.7. The same argument upstairs uses $\beta_X^2=0$ and $\int_X\beta_X\wedge\Omega>0$ for any Kähler $\Omega$. The latter is positive on the open isomorphism locus.

For this nef class, define the usual intersection numerical dimension as the largest $k$ with $\int_M\beta^k\wedge\omega_M^{2-k}>0$. It is exactly one. Replacing $\omega_M$ by any Kähler form cannot change that integer: positive forms on a compact manifold are bounded above and below by positive multiples of one another. It is also one on $X$ by the same smooth computations. For the actual nonpluripolar products of the smooth current, locality gives $\langle\beta\rangle_{\rm NP}=\beta$ and $\langle\beta^2\rangle_{\rm NP}=\beta^2=0$, and similarly upstairs. Their largest nonzero cohomology degree is one because of the positive pairing. Finally, the constant potential sequence $0,0,\ldots$ satisfies Cao Definition 2: convergence and positivity are immediate, singularity order is unchanged, and all multiplier ideals in its condition (iii) are $\mathcal O$. Definition 3 therefore computes $[\beta]$ and $[\beta^2]=0$. These three integers agree by these computations for this smooth example only. No equality is asserted for general fixed P-\allowbreak{}envelope data, actual products, Cao products or class numerical dimension.

Let $Y$ be any compact smooth curve on $X$. Its integration current pushes forward under $\pi$ to a closed positive $(1,1)$-\allowbreak{}current on $M$: this follows directly by testing against pullbacks of forms, by holomorphicity and Stokes on $Y$. Its cohomology class is integral because its homology class is the pushforward of the integral fundamental class of $Y$. Averaging the current under translations, just as for the smooth forms in Section 3, gives its constant $(1,1)$-\allowbreak{}representative. The choice of $M$ forces this integral class to be zero. But if $\pi|_Y$ were nonconstant, its derivative would be nonzero on a nonempty open part of $Y$, and $\int_Y\pi^*\omega_M>0$, contradicting that zero class. Thus $\pi(Y)$ is a point and

\[\beta_X|_Y=0,\qquad \pi^*\alpha\cdot Y=0.\]

So no such curve is a positive-\allowbreak{}mass one-\allowbreak{}dimensional complex carrier for this class. There is not even a globally smooth complete flag on $M$. The quantification over models is nonvacuous: blowing up any point of $M$ supplies an exceptional $\mathbb P^1$ and a flag centered at any of its points. Such a blowup is Kähler; explicitly the usual cutoff logarithm in a coordinate ball, with sufficiently small coefficient, pulls back after exceptional-\allowbreak{}divisor subtraction to a smooth Kähler form, exactly as reconstructed in the appendix's local calculation. No existence of a positive-\allowbreak{}volume subvariety has been assumed.

The null foliation does not provide a Hausdorff transverse quotient of positive dimension. Indeed any continuous map to a Hausdorff space which is constant on leaves is constant on all of $M$: the fiber over its value on one dense leaf is closed and contains that leaf. A measurable foliation space, local transverse geometry, or another enriched carrier is not excluded by this fact; none has yet been furnished for arbitrary prescribed P-\allowbreak{}data.

\subsection{6. Uniform collapse of both valuation coordinates}

Fix once and for all $X$, $Y$, $q$ and $\Omega$ as in the theorem. The class $\xi_\varepsilon=\pi^*\alpha+\varepsilon[\Omega]$ is Kähler, with smooth representative $\beta_X+\varepsilon\Omega$. Its DRWNXZ generating currents $G$ are gentle analytic Kähler currents. Write

\[t=\nu(G,Y),\qquad S=G-t[Y]\geq0,\qquad a=\nu(S|_Y,q).\]

Here $t$ is rational; published divisor subtraction applies. The residual is analytic along the whole flag curve, has generic Lelong number zero there, and so its restriction exists as a positive current on the compact curve, as in the source's recursive valuation definition. The valuation is $(t,a)$.

\textbf{First coordinate.} The trace mass $\int_XG\wedge\Omega=\int_X(\beta_X+\varepsilon\Omega)\wedge\Omega$ is bounded for $0<\varepsilon\leq1$. From $G\geq t[Y]$ follows a uniform upper bound on $t$, obtained by dividing that mass by the positive number $\int_Y\Omega$. In fact the supremum $f(\varepsilon)$ of these $t$ tends to zero. Otherwise there are $\varepsilon_i\downarrow0$ and generating currents $G_i$ with $t_i\geq c>0$. Their bounded trace masses give a weakly convergent subsequence of positive currents.

The precise compactness argument needed here is short. In finitely many coordinate charts positivity makes current coefficients Radon measures, whose total variations are bounded by a constant times their diagonal trace mass (the off-\allowbreak{}diagonal bound follows from positivity of each $2\times2$ coefficient matrix). The fixed Kähler trace bounds all these local masses on smaller charts. Weak compactness of finite measures and a finite-\allowbreak{}cover diagonal selection produce a current subsequence. Positivity, type and distributional closedness survive weak convergence. Pairing with each smooth closed complementary-\allowbreak{}degree form passes to the limit, so the limit represents $\pi^*\alpha$. The global potential representation and Section 4 force that limit to be the smooth current $\beta_X$.

After a further subsequence $t_i\to t\geq c$, since the $t_i$ are bounded. The inequalities $G_i-t_i[Y]\geq0$ pass to the weak limit and would give $\beta_X\geq t[Y]$. Wedge with $\Omega$ to obtain domination of trace Radon measures. The smooth left measure gives zero mass to $Y$; the right gives $t\int_Y\Omega>0$. This contradiction proves $f(\varepsilon)\to0$. It uses no unjustified interchange of generic Lelong numbers and weak limits.

\textbf{Second coordinate.} On a compact curve a positive $(1,1)$-\allowbreak{}current is a positive measure. In the present normalization its Lelong number at $q$ is its atom there, hence is at most its total degree. Put $b=\int_Y\Omega>0$ and $h=Y^2$. Cohomology of restriction and $\pi^*\alpha\cdot Y=0$ give

\[0\leq a\leq\int_Y S|_Y=(\xi_\varepsilon-t[Y])\cdot Y=\varepsilon b-th\leq\varepsilon b+|h|f(\varepsilon).\]

We need no assumption on the sign of $h$. All generating valuations, and hence their Euclidean closure, lie in

\[[0,f(\varepsilon)]\times[0,\varepsilon b+|h|f(\varepsilon)].\]

The smooth representative $\beta_X+\varepsilon\Omega$ is a gentle analytic Kähler current of valuation $(0,0)$. Thus the body contains the origin. Its Hausdorff distance to the origin is bounded by

\[\big(f(\varepsilon)^2+(\varepsilon b+|h|f(\varepsilon))^2\big)^{1/2}\longrightarrow0.\]

This proves both the Hausdorff limit and the intersection assertion. The proof is independent of the choice of $q$; it bounds the second coordinate by the entire curve degree.

At positive levels the smooth reference has potential zero. Every quasi-\allowbreak{}psh potential on compact $X$ is bounded above, so the inequality $v\leq0+C_v$ is automatic. Consequently every generating current satisfies the ordinary more-\allowbreak{}singular constraint, and the partial body is exactly the full body. All reference hypotheses of Definition 5.1 are satisfied. The fixed original P-\allowbreak{}datum is therefore present in the actual recipe, although it is the smooth one.

For a path $\rho_j$ of smooth closed Kähler forms tending uniformly to zero, repeat the compactness argument with $\beta_X+\rho_j$: their classes and masses converge to those of $\beta_X$. The supremum of the first coordinates tends to zero by the same contradiction. Replace $\varepsilon b$ in the second bound by $\int_Y\rho_j\to0$. The smooth reference again supplies the origin. This proves precisely the claimed path independence. Arbitrary nonpositive perturbations, moving models, moving flags or coordinate rescalings are not covered by this proof.

\subsection{7. What a rescaling would have to add}

There is genuine positive lower-\allowbreak{}degree data, but its scalar pairing needs a reference. For the fixed model and Kähler form put

\[A_\Omega=\int_X\beta_X\wedge\Omega>0,\qquad B_\Omega=\int_X\Omega^2>0.\]

The Kähler-\allowbreak{}class ambient volume (DRWNXZ Definition 2.12, LF 494–502) and the published surface volume formula give the exact positive-\allowbreak{}level area

\[2\operatorname{Area}\Delta_{Y,q}(\xi_\varepsilon)=\int_X(\beta_X+\varepsilon\Omega)^2=2\varepsilon A_\Omega+\varepsilon^2B_\Omega.\]

For a Kähler class, the equality of class volume with its smooth self-\allowbreak{}intersection can also be checked directly in that definition. Resolve a gentle analytic Kähler current as a smooth semipositive residual plus an effective divisor, using the third condition of gentleness. Positivity of the residual off the divisor implies its semipositivity everywhere by smoothness. Its smooth square has the same integral as the original current's nonpluripolar square: on the open isomorphism locus outside the divisor they agree by locality, and the smooth measure upstairs and the nonpluripolar measure downstairs give no mass to the omitted analytic sets. On a surface the difference between the square of the pulled-\allowbreak{}back smooth Kähler class and the residual square is the effective divisor paired with the sum of those two semipositive classes, hence nonnegative. A smooth representative attains the self-\allowbreak{}intersection. This proves the Kähler volume calculation for the source's supremum over gentle currents. No zero-\allowbreak{}mass normalization is involved.

Thus $\operatorname{Area}(\varepsilon^{-1/2}\Delta_{Y,q}(\xi_\varepsilon))\to A_\Omega$. This is only an area identity; it gives no boundedness, Hausdorff convergence, limiting shape, or one-\allowbreak{}dimensional carrier for those rescaled bodies. If that isotropically rescaled family had a compact Hausdorff limit of positive area, it would be two-\allowbreak{}dimensional. Furthermore it cannot have the same such limit for $\Omega$ and $c\Omega$ with $c\ne1$. Indeed

\[\varepsilon^{-1/2}\Delta(\pi^*\alpha+\varepsilon[c\Omega])=\sqrt c\,(c\varepsilon)^{-1/2}\Delta(\pi^*\alpha+c\varepsilon[\Omega]).\]

Whenever the right-\allowbreak{}hand rescaled limit exists it is $\sqrt c$ times the limit for $\Omega$, and its area is $cA_\Omega$. The area functional is continuous for compact convex sets in the Hausdorff topology: if the limit has interior, its boundary has Lebesgue measure zero and membership indicators converge almost everywhere in a common bounded box; if the limit has empty interior, its small neighborhoods have area tending to zero. This justifies the conditional area comparison.

One could choose a positive reference pairing $A_\Omega$ or rescale $\Omega$ to make that pairing one. This divides by a positive lower-\allowbreak{}degree number, not by zero top mass. It does not prove independence of different forms with the same pairing, and it does not select an anisotropic collapse, a geometric one-\allowbreak{}dimensional valuation, or a construction retaining general P-\allowbreak{}types. A numerical interval of length $A_\Omega$ alone is not an Okounkov realization. We do not declare one.

\subsection{8. The exact no-\allowbreak{}go conclusion and the continuing original task}

The following four requirements are incompatible:

\begin{enumerate}
\item[1.] The domain includes all zero-\allowbreak{}volume compact Kähler classes and prescribed ordinary P-\allowbreak{}types, in particular this $\alpha$ with P-\allowbreak{}type zero.
\item[2.] A chosen smooth Kähler modification and globally smooth complete flag give the usual recursive analytic-\allowbreak{}current valuations, with exactly the usual positive-\allowbreak{}level partial bodies subject to the fixed ordinary singularity constraint.
\item[3.] The zero-\allowbreak{}volume body is their \textbf{unrescaled} perturbative intersection, or their unrescaled Hausdorff limit along the positive paths proved above.
\item[4.] For this smooth nef datum its affine dimension equals one of the three numerical/\allowbreak{}product degrees explicitly computed in Section 5.

\end{enumerate}
Section 6 forces affine dimension zero for every permissible fixed modification and flag, while Section 5 computes each named degree to be one. This is a counterexample to those axioms together. It is stronger than choosing one bad flag, and it does not challenge Xia's projective conjecture or CHPW's positive-\allowbreak{}volume-\allowbreak{}subvariety theorem.

This reconstruction validates the proposed torus obstruction with those exact axioms. It does \textbf{not} prove impossibility for every richer geometric construction in original P15.3. No general construction is claimed. In particular the assignment's unchanged general obligations remain:

\begin{itemize}
\item For arbitrary prescribed ordinary P-\allowbreak{}envelope data, construct a genuine valuation or geometric realization, and prove its nonemptiness and compactness without replacing the data by I-\allowbreak{}equivalence. The retained dense-\allowbreak{}pole example in the appendix proves that direct analytic domination can have no competitor even at positive perturbation levels.
\item Specify an actual dimensional carrier and numerical invariant for that general construction. Fixed P-\allowbreak{}envelope products, representative nonpluripolar products, Cao current products, and class numerical dimension have not been identified in general. The product example in the appendix identifies them only where it explicitly computes them.
\item Handle spaces without smooth flags or positive-\allowbreak{}volume complex subvarieties. The torus proof shows that merely passing to a smooth model and then choosing a different usual smooth flag does not repair the numerical carrier in this example. It rules out neither nonstandard valuations nor a measured or other enriched geometric carrier.
\item State all reference scales and prove every claimed perturbation and model-\allowbreak{}choice independence. The positive-\allowbreak{}path limit in Section 6 is independent but has the wrong dimension for the named degrees. Section 7 computes a reference-\allowbreak{}sensitive rescaled area without supplying a limiting body. General rescaling and carrier independence remain unproved.

\end{itemize}
Several possible routes have therefore been tested only to the following extent. The ordinary analytic-\allowbreak{}domination route is defeated by the explicit dense-\allowbreak{}pole type. Keeping the class-\allowbreak{}level unrescaled valuation limit is defeated numerically by the torus. Passing to a Hausdorff null-\allowbreak{}leaf quotient fails for that torus. Isotropic rescaling has the reference and dimension issues proved in Section 7. These are limitations of those routes, not an exhaustive classification of all possible constructions. A general nonanalytic restriction/\allowbreak{}valuation theory, a non-\allowbreak{}Hausdorff or measured transverse carrier, or an appropriate anisotropic construction could evade the stated axioms; no proof of its existence or universal impossibility is supplied. There is no inference of impossibility from a failed search.

The preserved restricted theorem below is not promoted to original closure. P15.1 and P15.2 remain outside this assignment. There is no energy or Monge–Ampère operator construction and no normalization of a zero quantity to a probability. The final status is partial and \texttt{candidate=\allowbreak{}false}.

\subsection{9. Documentary evidence and rerun commands}

\texttt{input.json} and \texttt{input-\allowbreak{}original.json} are byte-\allowbreak{}identical. \texttt{assignment-\allowbreak{}original.txt} is the exact UTF-\allowbreak{}8 decoded assignment without an added newline. \texttt{task-\allowbreak{}original.json} is the decoded task serialized with Unicode, indentation and a terminal LF. Their derivations and original hashes are in \texttt{input-\allowbreak{}preservation.json}; they are never rewritten by the build scripts.

Run \texttt{python3 prepare\_\allowbreak{}sources.py} to copy and hash the four explicitly allowed caches and regenerate the literal LF extracts. Run \texttt{python3 assemble\_\allowbreak{}argument.py} to append the supplied historical proof with the two wording corrections (one correction of existing text and one source clarification), and generate the source/\allowbreak{}application ledger. Run \texttt{python3 build\_\allowbreak{}paper.py} to generate the standalone \texttt{P15.3-\allowbreak{}paper.tex} from the entire argument using Python and compile it with the existing XeLaTeX executable twice. Each actual build retains separate stdout, stderr and TeX logs. Run \texttt{python3 verify\_\allowbreak{}artifacts.py} to check input preservation, source extracts, ledger and paper consistency, and the exact status schema. OS and existing tool versions are recorded in \texttt{environment.json}; actual commands and exit codes are in the build manifest. No software is installed, no other document is opened or edited, no worker or other model is invoked, and no resource-\allowbreak{}heavy computation is attempted.

These are documentary checks. Hash matches, finite LF counts, successful compilation and program exit codes have no mathematical certification force. The new result rests on the analytic arguments in Sections 3–8 and the explicitly retained published background, not on computations. No numerical experiment or independent mathematical audit is claimed.

\section{Appendix: complete retained restricted proof}

The following is the supplied mathematical proof, with exactly the two harmless wording corrections recorded in the ledger. Its old documentary section is superseded by Section 9 above. The supplied source-\allowbreak{}aware PASS is retained in the frozen input as a review of this restricted scope only; it supplies no certification of the new torus theorem or of original closure.

\subsection{P15.3: a proved divisorial partial construction, with the original obligations left open}

This is an unverified mathematical draft. It proves a restricted construction, including a nonprojective transcendental example with class volume zero. It does \textbf{not} solve unchanged original P15.3. In particular, it gives no construction for arbitrary prescribed nonanalytic P-\allowbreak{}envelope data and no general numerical-\allowbreak{}dimension theorem. Compilation, source hashes, and successful execution are documentary checks, not mathematical verification.

\subsubsection{1. Sources and distinctions}

Only the four supplied public primary caches have been used. Their exact versions are DRWNXZ, arXiv:2309.07584v3 (5 January 2026); Xia, arXiv:2112.04290v2 (13 April 2024); Cao, arXiv:1210.5692v1 (21 October 2012); CHPW, arXiv:1508.03922v5 (arXiv header 5 February 2018, internal manuscript date 24 September 2018). Here and below source line numbers mean one-\allowbreak{}based indexing of \texttt{text.split('\textbackslash{}n')}, preserving form feeds and the trailing empty line. Literal statements and relevant proof bodies, their inclusive locations, and application checks are in \texttt{source-\allowbreak{}extracts/\allowbreak{}source-\allowbreak{}provenance.json}. Complete byte-\allowbreak{}identical supplied text extracts are also preserved there.

DRWNXZ Definition 3.2 and the pseudoeffective definition (lines 561–691) use successive generic Lelong numbers of analytic currents along a \textbf{globally smooth connected complete flag}. For a pseudoeffective class they take the intersection of the bodies of Kähler perturbations. Theorem 7.6 and Corollary 7.7 (lines 1389–1415) prove ambient volume equals class volume divided by n!. The last proof proves zero ambient volume by comparison with perturbations. It proves no lower-\allowbreak{}dimensional numerical equality and contains no prescribed-\allowbreak{}potential input. Their Definition 5.1 and Example 5.3 (lines 852–894) do impose bounded singularity order on analytic Kähler currents. That analytic partial construction, rather than just their class body, is used below.

Xia's sections/\allowbreak{}multiplier ideals, equation (1.2), and Proposition 5.8 (lines 73–99 and 1181–1185) make his partial body invariant under I-\allowbreak{}equivalence. His definitions distinguish bounded type, ordinary P-\allowbreak{}envelope, the mass-\allowbreak{}sensitive envelope C, and the I-\allowbreak{}envelope (lines 371–429). In particular, at zero actual mass Proposition 2.10 does not identify C with P. Theorem 2.19 (lines 490–494) computes the section volume as the actual mass of the \textbf{I-\allowbreak{}envelope}, divided by n!, not as an unconditional equality with the mass of the input potential. Positive class volume, positive ideal-\allowbreak{}model volume, and positive actual potential mass are distinct hypotheses. An input in a big class can have zero actual mass; an I-\allowbreak{}invariant body does not, merely from its definition, recover its ordinary P-\allowbreak{}envelope or actual mass.

Xia section 5.6 (lines 1623–1637) defines a limit using rational positive perturbations by an ample divisor and explicitly retains I-\allowbreak{}invariance. Conjecture 5.29 states \texttt{dim Delta(theta,phi) =\allowbreak{} nd(theta,phi)}; its next sentence refers to Cao Definition 4, and the following sentence expresses an expectation about a proof. It is not a proved theorem. Xia's source domain is projective and line-\allowbreak{}bundle based. It is not invoked for the transcendental example below.

Xia's section 6.4 (lines 1824–1840) leaves independence of the Chebyshev transform from its auxiliary more-\allowbreak{}singular Kähler-\allowbreak{}current potential eta open; the auxiliary datum is not an arbitrary positive form.

Cao Definition 2, Proposition 2.2 with its proof, and Definition 3 (lines 162–279) define a cohomological product through the absolutely continuous parts of quasi-\allowbreak{}equisingular analytic approximants. Definition 4 (lines 284–289) defines the numerical dimension of a line bundle with metric from that product. This is not the definition of the actual nonpluripolar product of the given current. We will denote by d\_\allowbreak{}Cao(T) the maximum nonzero degree of Cao's \textbf{current} product of Definition 3, which is stated for general positive currents; we do not pretend a genuinely transcendental class is a line bundle. Equality with actual products will be proved for the specific analytic type below.

CHPW Theorem B/\allowbreak{}Corollary 3.22 (lines 173–185 and 1071–1085) requires a smooth projective variety, a pseudoeffective real divisor, and a flag \textbf{containing a positive volume subvariety} of dimension equal to the numerical Iitaka dimension. The lower-\allowbreak{}dimensional volume is augmented restricted volume along that subvariety. The relevant proof bodies are Theorem 3.20 and Proposition 3.21 (lines 919–1070). Remark 3.24 (lines 1090–1096) expressly allows strict dimension inequality for other flags. None of this theorem is applied to our nonprojective example. Its dimension is established by an explicit product-\allowbreak{}body identity.

These are precise source limitations, not an impossibility proof for original P15.3.

\subsubsection{2. A restricted analytic construction}

Use dd\textasciicircum{}\allowbreak{}c =\allowbreak{} i partial-\allowbreak{}barpartial/\allowbreak{}(2 pi), so dd\textasciicircum{}\allowbreak{}c log|z|\textasciicircum{}\allowbreak{}2 is the unit divisor current. Fix a compact Kähler n-\allowbreak{}fold M with a globally smooth connected complete flag Y. Let F be a smooth prime divisor, a a positive rational number, and D=\allowbreak{}aF. Fix a Hermitian metric on O(F), with curvature e, and put

\begin{Verbatim}[breaklines=true,fontsize=\small]
phi = a log |s_F|_h^2,  theta = beta + a e,
T = theta + dd^c phi = beta + [D],  alpha = [theta],
\end{Verbatim}

where beta is a smooth closed semipositive (1,1)-\allowbreak{}form. Assume vol(alpha)=\allowbreak{}0. The chosen flag (or its specified flag on a modification) is part of the data. No independence of flags or arbitrary modifications is asserted.

For a Kähler form Omega on M, define

\begin{Verbatim}[breaklines=true,fontsize=\small]
B_Y(alpha,phi) = intersection over epsilon>0 of
    Delta_Y(alpha+epsilon[Omega]; T+epsilon Omega).
\end{Verbatim}

At each positive epsilon the right side is the genuine partial current-\allowbreak{}valuation body of DRWNXZ Definition 5.1: the reference current is Kähler, has rational gentle analytic divisorial singularities, and admissible currents must be more singular than phi up to a bounded constant. Thus this is a geometric Okounkov construction, not a numerical shape decorated with a potential.

\textbf{Proposition.} This construction is nonempty, compact, and convex, and

\begin{Verbatim}[breaklines=true,fontsize=\small]
B_Y(alpha,phi) = nu_Y([D]) + Delta_Y([beta]).                 (2.1)
\end{Verbatim}

It is independent of Omega. It retains the ordinary P-\allowbreak{}envelope type of phi within this divisorial bounded-\allowbreak{}type domain. Its dimension is dim Delta\_\allowbreak{}Y([beta]); no general identification of this dimension with a numerical dimension is claimed.

\textbf{Proof of (2.1).} At positive epsilon, R\_\allowbreak{}epsilon=\allowbreak{}[D]+beta+epsilon Omega satisfies every hypothesis of DRWNXZ Example 5.3. For completeness, its inclusions are as follows. A residual Kähler current S with gentle analytic singularities in [beta]+epsilon[Omega] yields S+[D], which is admissible for R\_\allowbreak{}epsilon, and its valuation is nu(S)+nu([D]). Conversely, an admissible analytic Kähler current G has generic divisorial multiplicity at least a on F. Divisor subtraction (DRWNXZ Lemma 2.5, including its analytic quotient proof) gives a positive analytic residual G-\allowbreak{}[D] in [beta]+epsilon[Omega]. It is Kähler: apply the same subtraction to G-\allowbreak{}delta Omega for a sufficiently small delta>0. Its valuation belongs to the residual class body by Definition 3.2/\allowbreak{}Lemma 3.1. Valuation additivity gives the converse inclusion after closure. Therefore

\begin{Verbatim}[breaklines=true,fontsize=\small]
Delta_Y(alpha+epsilon[Omega]; R_epsilon)
    = nu_Y([D]) + Delta_Y([beta]+epsilon[Omega]).             (2.2)
\end{Verbatim}

No nonanalytic regularization that supposedly preserves P-\allowbreak{}type is used in this step. Taking intersections and the actual pseudoeffective definition gives (2.1). Compactness and convexity follow from the published class-\allowbreak{}body construction. Nonemptiness also follows directly: the smooth residual beta+epsilon Omega gives nu([D]) in every body.

\textbf{Perturbation independence, with its exact scope.} Adding a smooth positive form to an analytic current does not change any valuation coordinate. Consequently, if xi and xi+[rho] are big and rho is smooth positive, Delta\_\allowbreak{}Y(xi) is contained in Delta\_\allowbreak{}Y(xi+[rho]). For two Kähler forms Omega and Omega', choose constants 0<c<C such that c Omega < Omega' < C Omega pointwise. The residual bodies for these perturbations sandwich one another after rescaling epsilon; hence their intersections are equal. This proves independence without comparing zero ambient volumes. More generally, if rho\_\allowbreak{}j are smooth Kähler forms tending to zero in the norm of a fixed Omega, then

\begin{Verbatim}[breaklines=true,fontsize=\small]
Delta_Y([beta]) subset Delta_Y([beta]+[rho_j])
    subset Delta_Y([beta]+delta_j[Omega]),  delta_j -> 0.
\end{Verbatim}

The right-\allowbreak{}hand compact bodies are nested when delta decreases. They converge in Hausdorff distance to their intersection: otherwise a sequence of points at fixed distance from the intersection has a convergent subsequence in a fixed compact body, whose limit belongs to every member. The sandwich therefore proves convergence to the residual class body along these paths; translating by nu\_\allowbreak{}Y([D]) gives (2.1). We claim neither arbitrary nonpositive perturbations nor convergence when the positive perturbations do not tend to zero. Rational cofinal epsilon sequences give the same intersection.

\textbf{Retained P-\allowbreak{}data.} Write

\begin{Verbatim}[breaklines=true,fontsize=\small]
P_theta[phi] = (sup {v in PSH(M,theta): v<=0, v<=phi+C_v})^*.
\end{Verbatim}

We prove P\_\allowbreak{}theta[phi]=\allowbreak{}phi+O(1), rather than presuming the I-\allowbreak{}envelope equals this envelope. Take a competitor v. Its divisorial Lelong number on F is at least a; subtraction gives a beta-\allowbreak{}psh potential w=\allowbreak{}v-\allowbreak{}phi, initially off F and then extended as the potential of the positive residual current. This extension is upper semicontinuous and agrees off F. On a compact set away from F, v<=\allowbreak{}0 bounds w above uniformly by -\allowbreak{}phi, independently of C\_\allowbreak{}v. In a coordinate polydisc along F, use F=\allowbreak{}\{z\_\allowbreak{}1=\allowbreak{}0\}, phi=\allowbreak{}a log|z\_\allowbreak{}1|\textasciicircum{}\allowbreak{}2+b with b smooth, and a local smooth potential r for beta. On |z\_\allowbreak{}1|=\allowbreak{}r\_\allowbreak{}0 and z' in a smaller polydisc, w+r is uniformly bounded above by -\allowbreak{}phi+r. Its restriction to each transverse disc is subharmonic or identically minus infinity. The maximum principle gives the same bound inside the disc. A finite covering of F, together with the compact complement, gives w<=\allowbreak{}K independent of v. Thus every competitor satisfies v<=\allowbreak{}phi+K, and this inequality persists after upper semicontinuous regularization, using the local logarithmic bound near F. The competitor phi-\allowbreak{}sup\_\allowbreak{}M phi gives the reverse bounded inequality. This proves the assertion.

If two inputs differ by a bounded function, their P-\allowbreak{}envelope competitor sets coincide. In particular P\_\allowbreak{}theta[P\_\allowbreak{}theta[phi]]=\allowbreak{}P\_\allowbreak{}theta[phi] here. Altering the Hermitian metric on O(F) changes phi by a smooth bounded function and does not change the body or the P-\allowbreak{}data. This is retention of the analytic P-\allowbreak{}envelope type through the actual singularity constraint; it does not assert that one body reconstructs arbitrary singularity types.

\textbf{Products in this restricted domain.} On M minus F, the current T is the smooth beta. Actual nonpluripolar products are local there and put no mass on F, so <T\textasciicircum{}\allowbreak{}k>\_\allowbreak{}NP=\allowbreak{}beta\textasciicircum{}\allowbreak{}k. The constant analytic approximation phi\_\allowbreak{}j=\allowbreak{}phi satisfies all of Cao Definition 2: L1 convergence and positivity are immediate, singularity order is unchanged, and the required multiplier-\allowbreak{}ideal inclusion for (1+delta)m follows after normalizing phi<=\allowbreak{}0. The absolutely continuous part of T is beta. Cao Definition 3 therefore gives [beta\textasciicircum{}\allowbreak{}k]. Its equality with [<T\textasciicircum{}\allowbreak{}k>\_\allowbreak{}NP] has now been proved for this domain, without making a general identification.

The cohomological data also persist for a positive current of the same bounded singularity type. After subtracting [D], its residual potential b is globally bounded relative to beta. The bounded Bedford-\allowbreak{}-\allowbreak{}Taylor products have cohomology [beta\textasciicircum{}\allowbreak{}k]: the telescoping identity

\begin{Verbatim}[breaklines=true,fontsize=\small]
(beta+dd^c b)^k - beta^k
  = dd^c( b sum_{j=0}^{k-1} (beta+dd^c b)^j wedge beta^{k-1-j})
\end{Verbatim}

is valid for bounded potentials by the defining local Bedford-\allowbreak{}-\allowbreak{}Taylor operation and its multiplication by closed products. Off F these are the nonpluripolar products, and bounded products charge no pluripolar set. Thus their extension across F is the actual nonpluripolar product. This proves cohomological invariance for the retained type. It is not equality of the different representative measures.

\subsubsection{3. A genuinely transcendental zero-\allowbreak{}class-\allowbreak{}volume example with a proved dimension theorem}

Here is an explicit admissible domain where (2.1) has proved numerical content. Let S be a compact Kähler surface with a Kähler class kappa outside the real span of divisor classes; fix a Kähler form in it. Such input may, for example, be a two-\allowbreak{}dimensional complex torus with no nonzero integral (1,1)-\allowbreak{}classes. Existence of such tori follows by choosing a linear complex structure on R\textasciicircum{}\allowbreak{}4 outside the countably many proper real-\allowbreak{}analytic loci on which a fixed nonzero integral alternating two-\allowbreak{}form is of type (1,1); each locus is proper since a complex structure can be varied to fail that invariance condition. Baire's theorem supplies the choice. A constant Hermitian form makes the quotient R\textasciicircum{}\allowbreak{}4/\allowbreak{}Z\textasciicircum{}\allowbreak{}4 Kähler. Only existence, not any computation on a selected torus, is being asserted.

Fix p in S. Put pi: S\_\allowbreak{}hat=\allowbreak{}Bl\_\allowbreak{}p S -\allowbreak{}> S, with exceptional E. Choose coordinates near p and a smooth cutoff chi supported in that chart, equal to 1 near p. Let

\begin{Verbatim}[breaklines=true,fontsize=\small]
f=chi log(|z_1|^2+|z_2|^2), u_a=a f.
\end{Verbatim}

For every sufficiently small a>0, kappa\_\allowbreak{}form+dd\textasciicircum{}\allowbreak{}c u\_\allowbreak{}a is positive. Indeed dd\textasciicircum{}\allowbreak{}c log|z|\textasciicircum{}\allowbreak{}2 is semipositive away from p in the region chi=\allowbreak{}1, and the negative part of dd\textasciicircum{}\allowbreak{}c f on the compact cutoff annulus is bounded by a fixed multiple of the Kähler form. Choose a positive rational a below the resulting bound. Its pullback has the decomposition

\begin{Verbatim}[breaklines=true,fontsize=\small]
pi^*(kappa_form+dd^c u_a)=beta_a+a[E],
[beta_a]=kappa_a=pi^*kappa-a[E],                            (3.1)
\end{Verbatim}

with beta\_\allowbreak{}a a \textbf{smooth Kähler form}. To see smoothness and strict positivity near E, in the chart z\_\allowbreak{}1=\allowbreak{}w, z\_\allowbreak{}2=\allowbreak{}wt the potential is a log|w|\textasciicircum{}\allowbreak{}2+a log(1+|t|\textasciicircum{}\allowbreak{}2). Subtraction of a[E] leaves pi\textasciicircum{}\allowbreak{}*kappa\_\allowbreak{}form+a dd\textasciicircum{}\allowbreak{}c log(1+|t|\textasciicircum{}\allowbreak{}2); on E the first form is positive in the normal direction and the second in the tangent direction. They are semipositive in a neighborhood with their kernels complementary. The other chart is identical. On the cutoff annulus small a gives strict positivity, and outside it beta\_\allowbreak{}a is the pullback of kappa\_\allowbreak{}form where pi is an isomorphism. This proves (3.1), not just a nef or unproved Seshadri assumption.

Let C be any connected compact Riemann surface, c\_\allowbreak{}0 a point, and q a point of E. Start on X\_\allowbreak{}0=\allowbreak{}S x C with class alpha\_\allowbreak{}0=\allowbreak{}p\_\allowbreak{}S\textasciicircum{}\allowbreak{}*kappa and potential phi\_\allowbreak{}0=\allowbreak{}p\_\allowbreak{}S\textasciicircum{}\allowbreak{}*u\_\allowbreak{}a. Use the specified modification

\begin{Verbatim}[breaklines=true,fontsize=\small]
Pi=pi x id_C: X=S_hat x C -> X_0.
\end{Verbatim}

Its center \{p\} x C is smooth. On X put alpha=\allowbreak{}Pi\textasciicircum{}\allowbreak{}*alpha\_\allowbreak{}0, theta=\allowbreak{}p\_\allowbreak{}S\_\allowbreak{}hat\textasciicircum{}\allowbreak{}*pi\textasciicircum{}\allowbreak{}*kappa\_\allowbreak{}form, phi\_\allowbreak{}a=\allowbreak{}Pi\textasciicircum{}\allowbreak{}*phi\_\allowbreak{}0, beta=\allowbreak{}p\_\allowbreak{}S\_\allowbreak{}hat\textasciicircum{}\allowbreak{}*beta\_\allowbreak{}a, and D=\allowbreak{}a(E x C). Then theta+dd\textasciicircum{}\allowbreak{}c phi\_\allowbreak{}a=\allowbreak{}beta+[D]. The globally smooth connected complete flag is

\begin{Verbatim}[breaklines=true,fontsize=\small]
X superset Y_1=S_hat x {c_0}
  superset Y_2=E x {c_0} superset Y_3={(q,c_0)}.             (3.2)
\end{Verbatim}

It is a current-\allowbreak{}valuation flag. We do not require the meromorphic function field of a nonprojective manifold to admit a surjective rank-\allowbreak{}three valuation.

The example is genuinely transcendental: if alpha\_\allowbreak{}0 belonged to the real span of divisor classes, restricting to S x \{c\} would put kappa in that span. The same argument after pullback uses restriction followed by pi\_\allowbreak{}*; divisor classes push to divisor classes or zero. Thus alpha is also outside that span. X\_\allowbreak{}0 is nonprojective for the indicated torus example. The modification supplies the required flag even when a suitable flag was unavailable on the original manifold.

We need a psh pullback statement at zero class volume, whereas the setup of DRWNXZ section 2 initially fixes a big class. Here is its proof with no bigness hypothesis. If v is Pi\textasciicircum{}\allowbreak{}*theta\_\allowbreak{}0-\allowbreak{}psh on X, then away from the exceptional image, w=\allowbreak{}v composed with Pi\textasciicircum{}\allowbreak{}\{-\allowbreak{}1\} is theta\_\allowbreak{}0-\allowbreak{}psh and is bounded above, because v is quasi-\allowbreak{}psh on a compact space. Locally add a smooth potential of theta\_\allowbreak{}0. A psh function bounded above extends over an analytic set by its upper semicontinuous extension (the removable-\allowbreak{}singularity argument used in the proof of DRWNXZ Lemma 2.1). Its pullback and v are psh potentials agreeing outside an analytic set, hence almost everywhere; canonical psh representatives agreeing almost everywhere agree everywhere by the submean regularization. This proves surjectivity. Injectivity and pointwise order preservation follow from surjectivity of Pi. This reproduces the cited proof and verifies the extension to the present nonbig class rather than assuming a big-\allowbreak{}class theorem applies.

A surjective pullback preserves both inequalities defining a bounded type and v<=\allowbreak{}0. Ordinary envelopes are least psh upper bounds of their competitor families, so the proved order isomorphism also commutes with P-\allowbreak{}envelopes for the pulled-\allowbreak{}back theta and zero obstacle. The retained P-\allowbreak{}type is consequently the pullback of the prescribed original point-\allowbreak{}logarithmic P-\allowbreak{}type. No independence of unrelated modifications or flags is asserted.

\textbf{Theorem (restricted product case).} For these data,

\begin{Verbatim}[breaklines=true,fontsize=\small]
B_Y(alpha,phi_a)
  = {(0,a+x,y): (x,y) in Delta_(E,q)(kappa_a)}.              (3.3)
\end{Verbatim}

Its affine carrier is \{x\_\allowbreak{}1=\allowbreak{}0\}, with the standard (x\_\allowbreak{}2,x\_\allowbreak{}3) coordinate area. Its intrinsic dimension is exactly 2, and

\begin{Verbatim}[breaklines=true,fontsize=\small]
2! Area(B_Y)=integral_(S_hat) beta_a^2
             = integral_S kappa_form^2 - a^2 > 0.         (3.4)
\end{Verbatim}

The ambient three-\allowbreak{}dimensional volume, class volume, and actual top nonpluripolar mass are all zero. The degree-\allowbreak{}two actual nonpluripolar cohomological data and Cao current product are [beta\textasciicircum{}\allowbreak{}2], nonzero; their largest nonzero degree is 2. In this example the nef class alpha itself also has numerical dimension 2. Equation (3.4) is an unnormalized lower-\allowbreak{}dimensional mass, not a probability.

\textbf{Class volume zero.} alpha has the smooth semipositive representative theta. It is not big: a Kähler current K in alpha would satisfy K>=\allowbreak{}delta Omega for some Kähler Omega, whereas

\begin{Verbatim}[breaklines=true,fontsize=\small]
integral_X K wedge theta^2 = integral_X theta^3 = 0,
integral_X Omega wedge theta^2 > 0.
\end{Verbatim}

This is a contradiction. The published definition assigns volume zero to a nonbig pseudoeffective class. The same proof works on X\_\allowbreak{}0. No implication from a particular potential's zero mass to class volume zero is used.

\textbf{The product-\allowbreak{}body identity supplying the dimension proof.} Choose a Kähler form eta on S\_\allowbreak{}hat and a positive form sigma on C of degree 1. Let Omega\_\allowbreak{}0=\allowbreak{}p\_\allowbreak{}S\_\allowbreak{}hat\textasciicircum{}\allowbreak{}*eta+p\_\allowbreak{}C\textasciicircum{}\allowbreak{}*sigma. We prove, for every epsilon>0,

\begin{Verbatim}[breaklines=true,fontsize=\small]
Delta_Y([beta]+epsilon[Omega_0])
   = [0,epsilon] x Delta_(E,q)(kappa_a+epsilon[eta]).       (3.5)
\end{Verbatim}

For the upper inclusion, take a positive analytic current G in this class, with first valuation t along Y\_\allowbreak{}1. Subtracting t[Y\_\allowbreak{}1] leaves a positive analytic current. Restrict it to a generic curve \{z\} x C avoiding the locus on which the analytic potential is identically minus infinity. Its degree is epsilon-\allowbreak{}t, since Y\_\allowbreak{}1 intersects that curve once; hence 0<=\allowbreak{}t<=\allowbreak{}epsilon. The restriction to Y\_\allowbreak{}1 after the first subtraction is positive, analytic near the flag point, and lies in kappa\_\allowbreak{}a+epsilon[eta]. (The normal bundle of a product fiber is trivial; the curve pullback contributes no surface class.) Its remaining two coordinates lie in the surface body by DRWNXZ Definition 3.2/\allowbreak{}Lemma 3.1. The upper inclusion follows after closure.

For the lower inclusion take a surface Kähler current R with gentle analytic singularities in kappa\_\allowbreak{}a+epsilon[eta]. For any rational 0<=\allowbreak{}t<epsilon, the product current

\begin{Verbatim}[breaklines=true,fontsize=\small]
p_S_hat^* R + p_C^*((epsilon-t)sigma+t[c_0])
\end{Verbatim}

is Kähler with rational gentle analytic singularities and has valuation (t,nu\_\allowbreak{}(E,q)(R)). Its positivity is strict in both product directions. Rational t are dense in [0,epsilon], including the endpoint in closure, and surface valuations of these R are dense in the surface body. This gives the full lower inclusion and proves (3.5).

Now kappa\_\allowbreak{}a is Kähler. Proposition 3.3(3) of DRWNXZ, applied to this \textbf{big surface class}, gives

\begin{Verbatim}[breaklines=true,fontsize=\small]
intersection_(epsilon>0) Delta_(E,q)(kappa_a+epsilon[eta])
  = Delta_(E,q)(kappa_a).
\end{Verbatim}

Intersecting (3.5) therefore gives Delta\_\allowbreak{}Y([beta])=\allowbreak{}\{0\} x Delta\_\allowbreak{}(E,q)(kappa\_\allowbreak{}a). The fixed divisor has valuation (0,a,0): it does not contain Y\_\allowbreak{}1, restricts to aE on Y\_\allowbreak{}1, and has zero residual after subtracting that divisor. Equation (2.1) proves (3.3). The ordinary surface class body has nonempty interior because kappa\_\allowbreak{}a is Kähler. Hence the dimension is 2; it is not inferred from Xia's conjecture or from zero ambient volume.

\textbf{Numerical mass and singularity sensitivity.} Apply DRWNXZ Theorem 7.6 to kappa\_\allowbreak{}a on the surface. Its volume is its Kähler self-\allowbreak{}intersection: for bounded truncations all top masses equal this intersection by Stokes, arbitrary nonpluripolar masses are at most that value, and a smooth representative attains it. Since pi\textasciicircum{}\allowbreak{}*kappa.E=\allowbreak{}0 and E\textasciicircum{}\allowbreak{}2=\allowbreak{}-\allowbreak{}1, expansion gives (3.4). These intersection facts follow respectively from restricting a pulled-\allowbreak{}back class to the contracted E, and from the normal bundle O\_\allowbreak{}E(-\allowbreak{}1). The actual nonpluripolar products of the chosen current are beta\textasciicircum{}\allowbreak{}k by Section 2, so the top degree 3 vanishes and degree 2 is nonzero. The constant analytic approximation computes exactly the same Cao cohomological products. For the nef class alpha the smooth representative theta likewise has largest nonzero intersection degree 2.

The surface geometric carrier is Y\_\allowbreak{}1 and its residual two-\allowbreak{}dimensional mass is integral beta\_\allowbreak{}a\textasciicircum{}\allowbreak{}2. Equivalently, with the explicitly named transverse scale integral\_\allowbreak{}C sigma=\allowbreak{}1,

\begin{Verbatim}[breaklines=true,fontsize=\small]
integral_X beta^2 wedge Omega_0 = integral_(S_hat) beta_a^2.
\end{Verbatim}

Changing eta or the degree-\allowbreak{}one sigma does not change this number. An arbitrary reference form of fiber degree m would instead give m times this number; there is no claim that that changed reference scale is irrelevant. No zero quantity is divided out.

The body contains (0,a,0), because the surface body contains (0,0) from a smooth Kähler representative; all its second coordinates are at least a. Thus min\_\allowbreak{}B x\_\allowbreak{}2=\allowbreak{}a. For fixed p, this recovers the rational logarithmic coefficient of the prescribed P-\allowbreak{}type. Distinct a give distinct bodies in the \textbf{same original class} alpha\_\allowbreak{}0; their lower mass is kappa\textasciicircum{}\allowbreak{}2-\allowbreak{}a\textasciicircum{}\allowbreak{}2. This demonstrates actual type sensitivity within the stated family, rather than appending a label to a class body. It does not show recovery of arbitrary supports or nonanalytic P-\allowbreak{}types from one flag.

All Kähler perturbation choices on the specified modified manifold, and the positive paths explicitly covered in Section 2, have the same limit. The proof uses the positive-\allowbreak{}epsilon valuation realization and compact sandwiches, not numerical volume equality at zero.

\subsubsection{4. An exact obstruction to the direct analytic-\allowbreak{}domination recipe}

The limitation to analytic divisorial input is substantive. There are prescribed P-\allowbreak{}envelope data at class volume zero for which \textbf{no} nontrivial competitor more singular than that envelope is analytic even near a single flag point, for any positive Kähler perturbation. This rules out the direct extension of the positive-\allowbreak{}epsilon competitor construction of Section 2 to all P-\allowbreak{}types; it does not rule out a different geometric construction.

Take M=\allowbreak{}P\textasciicircum{}\allowbreak{}1 x C for a compact connected curve C, theta=\allowbreak{}p\_\allowbreak{}1\textasciicircum{}\allowbreak{}*omega\_\allowbreak{}FS, with integral\_\allowbreak{}(P\textasciicircum{}\allowbreak{}1) omega\_\allowbreak{}FS=\allowbreak{}1. Its class is pseudoeffective and nonbig: a Kähler current in [theta], paired with theta, contradicts theta\textasciicircum{}\allowbreak{}2=\allowbreak{}0 and integral Omega wedge theta>0. Hence its class volume is zero. Choose distinct dense points p\_\allowbreak{}j in P\textasciicircum{}\allowbreak{}1 and a\_\allowbreak{}j=\allowbreak{}2\textasciicircum{}\allowbreak{}\{-\allowbreak{}j-\allowbreak{}2\}, j>=\allowbreak{}1, so A=\allowbreak{}sum a\_\allowbreak{}j=\allowbreak{}1/\allowbreak{}4. Let g\_\allowbreak{}j be the logarithm of the squared chordal distance to p\_\allowbreak{}j. It satisfies g\_\allowbreak{}j<=\allowbreak{}0, dd\textasciicircum{}\allowbreak{}c g\_\allowbreak{}j=\allowbreak{}[p\_\allowbreak{}j]-\allowbreak{}omega\_\allowbreak{}FS, has unit logarithmic pole at p\_\allowbreak{}j, and its L1 norm with respect to omega\_\allowbreak{}FS is independent of p\_\allowbreak{}j by unitary invariance. Define

\begin{Verbatim}[breaklines=true,fontsize=\small]
u=sum_(j>=1) a_j p_1^*g_j.
\end{Verbatim}

The partial sums are decreasing theta-\allowbreak{}psh functions with uniformly summable L1 norms. Thus their limit is theta-\allowbreak{}psh and not identically minus infinity; in particular it is a legitimate input. Its current is

\begin{Verbatim}[breaklines=true,fontsize=\small]
theta+dd^c u = (1-A)theta + sum_(j>=1) a_j[F_j],
F_j={p_j} x C.
\end{Verbatim}

The equality follows in distributions from L1 convergence; positivity is explicit. Its top nonpluripolar mass also vanishes: every bounded truncation max(u,-\allowbreak{}k) is pulled back from P\textasciicircum{}\allowbreak{}1, so its local top wedge on the surface is zero. This potential's zero mass is established separately from class volume zero.

Let psi=\allowbreak{}P\_\allowbreak{}theta[u], with the zero obstacle as in Section 2. The competitor u is bounded above by zero, so the envelope is nontrivial. For each j, u<=\allowbreak{}a\_\allowbreak{}j p\_\allowbreak{}1\textasciicircum{}\allowbreak{}*g\_\allowbreak{}j. Consequently every competitor for P\_\allowbreak{}theta[u] is a competitor for P\_\allowbreak{}theta[a\_\allowbreak{}j p\_\allowbreak{}1\textasciicircum{}\allowbreak{}*g\_\allowbreak{}j]. The latter is the single smooth-\allowbreak{}divisor situation of the envelope proof in Section 2, with residual (1-\allowbreak{}a\_\allowbreak{}j)theta. Thus, for a constant K\_\allowbreak{}j independent of the competitors,

\begin{Verbatim}[breaklines=true,fontsize=\small]
psi <= P_theta[a_j p_1^*g_j]
     <= a_j p_1^*g_j+K_j.
\end{Verbatim}

Upper semicontinuous regularization preserves this logarithmic upper bound by the same transverse maximum argument. In particular psi=\allowbreak{}-\allowbreak{}infinity along every F\_\allowbreak{}j. Their union is dense in M.

Now let v be any nontrivial quasi-\allowbreak{}psh potential satisfying v<=\allowbreak{}psi+C. It is minus infinity on every F\_\allowbreak{}j. If v had analytic singularities on any nonempty open neighborhood U, its minus-\allowbreak{}infinity locus in a smaller coordinate neighborhood would be the zero set of finitely many holomorphic functions, a proper closed analytic subset (or empty). It would contain the dense set U intersect union F\_\allowbreak{}j; hence all of that coordinate neighborhood. All generators would vanish identically, making v identically minus infinity there, contrary to a nontrivial quasi-\allowbreak{}psh potential's local integrability. This proves the impossibility.

The argument does not use the background class of v. It applies in particular to v in PSH(M,theta+epsilon Omega), for every Kähler Omega and epsilon>0. Even restricting to currents analytic only near a chosen flag center cannot make that competitor set nonempty. This defeats the direct recipe "value analytic currents more singular than the fixed P-\allowbreak{}envelope, then take perturbative intersections" for this explicit prescribed datum. It does not defeat the unchanged original demand if another genuine valuation/\allowbreak{}geometric realization can incorporate such data; none is supplied here. It is not an inference of impossibility from an unsuccessful search.

\subsubsection{5. What remains of unchanged original P15.3}

The original quantification includes arbitrary compact Kähler manifolds, pseudoeffective classes with zero class volume, and prescribed P-\allowbreak{}singularity data. The preceding proof handles a rational analytic divisorial type with smooth semipositive residual; its numerical equality is proved only in the product example. The following obligations remain unresolved:

\begin{enumerate}
\item[1.] Construct a genuine current-\allowbreak{}valuation or geometric Okounkov body for arbitrary prescribed nonanalytic P-\allowbreak{}envelope type, without replacing that type by I-\allowbreak{}equivalence. The analytic divisor-\allowbreak{}subtraction argument does not approximate such types while retaining them. Section 4 proves that simply keeping the same analytic-\allowbreak{}domination competitor rule can yield an empty set at every positive perturbation.
\item[2.] Establish nonemptiness, compactness, limiting behavior, and any P-\allowbreak{}preserving analytic approximation required by that construction. These do not follow from I-\allowbreak{}equisingular approximation.
\item[3.] Define and prove a general lower-\allowbreak{}dimensional numerical theorem for the retained type, with the relevant actual product or other numerical invariant precisely specified. Neither Cao's product nor Xia's conjectured dimension equality has been identified with such a theorem here.
\item[4.] Prove the exact flag hypotheses realizing that numerical dimension. An arbitrary flag cannot be assumed to do so. The product proof establishes only (3.2).
\item[5.] Specify a modification when a smooth complete flag is absent, and prove every claimed relation or independence under those choices. We retain one specified modification and flag; no general independence theorem has been proved.
\item[6.] Name any further reference scales or carriers in the general construction and prove every claimed perturbation independence. The positivity sandwiches here require a fixed divisor and smooth residual; they do not justify a general P-\allowbreak{}sensitive limit.

\end{enumerate}
There is no proof here that the original demand is impossible. The source gaps, restricted scope, and obstruction to one specified recipe are not such a proof. P15.1 and P15.2 are outside this assignment and no energy or Monge-\allowbreak{}-\allowbreak{}Ampère operator is constructed.


\section{Source/\allowbreak{}application ledger}

The machine-\allowbreak{}readable ledger is \texttt{source-\allowbreak{}application-\allowbreak{}ledger.json}. The full literal extracts and their inclusive LF coordinates are \texttt{source-\allowbreak{}extracts/\allowbreak{}source-\allowbreak{}provenance.json}.

\textbf{S1 — drwnxz, LF 226–263; 314–385; 561–691.} Smooth complete flags; rational recursive analytic-\allowbreak{}current valuations; big-\allowbreak{}class body closure and pseudoeffective intersection.

Application: New torus theorem Section 6 and retained divisorial/\allowbreak{}product construction. Hypotheses: X smooth connected compact Kähler surface; fixed globally smooth complete flag; positive perturbation representative beta\_\allowbreak{}X+rho is Kähler. Generators gentle analytic Kähler; rational divisorial coefficient. No recursive restriction of arbitrary nonanalytic potential is asserted. Status: Published lemma/\allowbreak{}definition/\allowbreak{}proofs retained. Uniform collapse is proved here, not supplied by these statements.

\textbf{S2 — drwnxz, LF 255–263.} Potential pullback/\allowbreak{}descent proof uses bounded-\allowbreak{}above psh removability and equality of canonical representatives off an analytic set.

Application: New Section 4; retained product modification. Hypotheses: Proper bimeromorphic morphism, smooth compact Kähler source and target; global upper bound. The source standing assumption is big, but this draft reconstructs its proof without bigness. Status: Nonbig extension proved in the draft; exact classical removability statement is explicitly retained from the published proof.

\textbf{S3 — drwnxz, LF 338–361; 852–894.} Partial bodies for gentle analytic Kähler references; divisorial translate identity and analytic realization.

Application: Smooth-\allowbreak{}reference partial bodies in Section 6; rational divisorial appendix. Hypotheses: Smooth reference at positive torus levels; every potential bounded above, so zero-\allowbreak{}type comparison automatic. Appendix D rational, residual smooth semipositive; epsilon makes residual Kähler. No general P-\allowbreak{}type replacement is used. Status: Partial identity is reproved in the appendix; no application to arbitrary dense-\allowbreak{}pole references.

\textbf{S4 — drwnxz, LF 494–502; 939–963; 1139–1415.} Volume zero for nonbig classes; class ambient volume formula and its induction proof.

Application: New Section 7 perturbative area; retained big surface mass/\allowbreak{}area calculation. Hypotheses: Formula applied only to Kähler/\allowbreak{}big surface positive levels or the Kähler residual surface in appendix. Nonbig volume is separately proved via a positive test pairing. No lower numerical dimension conclusion from zero ambient volume. Status: Published theorem used as background. Kähler self-\allowbreak{}intersection calculation separately proved here.

\textbf{S5 — xia, LF 67–105; 371–429; 476–494; 1158–1185; 1623–1637; 1824–1840.} Projective line-\allowbreak{}bundle multiplier-\allowbreak{}ideal bodies; ordinary, C, and I distinctions; section volume from I-\allowbreak{}envelope mass; zero-\allowbreak{}volume dimension conjecture; eta-\allowbreak{}dependent Chebyshev question.

Application: Source limitations only; no numerical result applied to nonprojective torus. Hypotheses: Positive ideal-\allowbreak{}model/\allowbreak{}section volume, positive actual input mass, and big class are distinct. Ordinary P-\allowbreak{}envelope is used as defined. Chebyshev auxiliary datum is a more-\allowbreak{}singular Kähler-\allowbreak{}current potential eta. Status: Definitions and literal statements/\allowbreak{}proof references retained; Conjecture 5.29 remains conjectural.

\textbf{S6 — cao, LF 162–279; 281–289; 660–721.} Global potential representation and quasi-\allowbreak{}equisingular current products; numerical dimension for line-\allowbreak{}bundle metrics; projective general-\allowbreak{}hypersurface reduction.

Application: Uniqueness of positive current from rigidity; constant approximations compute products in Sections 5 and appendix. Hypotheses: Currents on compact Kähler manifolds; constant analytic approximations verify all three conditions directly. Definition 3 applies to transcendental smooth currents. Definition 4 is not asserted for a nonexistent line bundle. General ample hypersurface argument not applied. Status: Published potential/\allowbreak{}product background explicitly retained; equality with actual products established only for the stated smooth/\allowbreak{}divisorial examples.

\textbf{S7 — chpw, LF 166–185; 530–577; 915–1096.} Numerical dimension and lower intrinsic volume for projective real divisors with a positive-\allowbreak{}volume-\allowbreak{}subvariety flag; other flags may have smaller dimension.

Application: Source limitation, not a theorem used to prove new collapse or old transcendental dimension. Hypotheses: Smooth projective variety, real divisor, flag containing positive volume subvariety. They are absent for the torus. All relevant proof bodies are extracted. Status: Published result retained without any claimed transcendental extension.

Contraction of model curves is proved by pushing forward their integral integration currents. The bounded-\allowbreak{}potential/\allowbreak{}nonpluripolar background in the supplied appendix is retained only in its stated scope; no fifth cache or book theorem was read. The torus existence, leaf density, rigidity, nonbig descent, weak compactness argument, valuation uniformity and reference-\allowbreak{}sensitive area assertions are proved in the main argument.

\end{document}
